%% sections/06-discussion.tex

\section{Discussion}
\label{sec:discussion}

\subsection{Player Styles as Design Tools}

The most practically significant finding is that all 7 styles are design-enabled rather
than personality-driven. A designer who wants sheet-deep-reader behavior can write
PC sheets with detailed grief-paragraphs, local-origin notes, and institutional training
backgrounds. A designer who wants craft-witness behavior can write PC sheets with named
daily practices. The design input is specific and learnable; the behavioral output is
reliable.

This reframes the role of TRPG character-sheet design in the research literature.
Character sheets are typically analyzed as mechanical artifacts (stat blocks, ability
scores, spell lists) or narrative artifacts (backstory, personality, alignment). The
marathon styles suggest a third function: character sheets are \textit{behavioral design
artifacts} that specify how a PC will engage with campaign events. The quality of this
behavioral specification --- whether it is specific, named, and connected to the
campaign's material --- determines what player styles the PC can exhibit.

\subsection{Cross-Campaign Archetype Independence}

The finding that all 7 styles fire across multiple campaign archetypes, including
Campaign 7 (no designed personal stakes, no grief paragraphs, no professional codes)
is the most analytically significant result. Campaign 7 produced craft-witness instances
(Vess Ironmark's structural-assessment habit), sheet-deep-reader instances (Brek's
regional knowledge producing terrain advantages), act-without-announcement instances
(Tam's road-companion behavior following Brek's wrong-fork decision), and npc-arc-completion
trigger instances (Elder Sorr, Cassel, Torven all firing by character logic).

Campaign 7 PCs were designed as strangers without grief hooks. The styles fired not
from grief but from situational professional habits, regional knowledge, and
relationship-driven behavioral change. This is the RV9 research validation: player
styles without ANY designed personal stakes are possible, provided the character
sheets are designed with sufficient specificity.

\subsection{The 50-Instance Sheet-Deep-Reader Question}

Sheet-deep-reader's 50+ instances across 7 campaigns suggests it is the most
fundamental of the seven styles. Every PC sheet with non-generic content (specific
origin, named relationships, institutional training) enables sheet-deep-reader behavior.
The style is less a distinct design pattern than the natural consequence of dense
PC-sheet design.

This implies that the 50-instance count is a lower bound, not a peak. In a corpus with
denser PC sheets or longer campaigns, sheet-deep-reader instances would be proportionally
higher. The style functions as a baseline measure of PC-sheet quality: the number of
sheet-deep-reader instances per session is a proxy for how much campaign-relevant
content the PC sheets contain.

\subsection{Limitations}

All instances were observed in AI-simulated sessions where the ``player'' has no
personality, emotional state, or social dynamic --- only the character sheet and
Playstyle Heuristics. Whether human players with the same PC sheets would exhibit
the same patterns is an empirical question. We expect that human players would produce
many of the same instances (the sheet design enables the behavior regardless of player
personality) and additional instances the sheet does not anticipate (human improvisational
capacity exceeds the AI's adherence to heuristics). Human play would also introduce
instances that are \textit{not} sheet-driven (table banter, player personality
bleeding through), which would increase measurement noise.
